\section{Fase 1. Evaluación del sistema original}

La inspección inicial orientó la intervención hacia los mecanismos de guiado y traslación. El contacto entre las piezas plásticas y los tubos se interpretó como un factor que contribuyó a la resistencia observada durante el movimiento. Esta interpretación se basó en la inspección del conjunto y no en una medición independiente de las fuerzas de fricción.

En X y Y, la desalineación y el avance desigual de los extremos del pórtico se relacionaron con las dificultades de desplazamiento de los carros. El comportamiento de tipo \textit{gantry racking} mostró que el movimiento del conjunto no fue uniforme. No se aisló experimentalmente la contribución de cada componente, por lo que la evaluación identificó condiciones de montaje asociadas con el problema, sin establecer una causa única.

En Z, la desalineación del motor, el mecanismo de desplazamiento y la garra indicó la necesidad de revisar la geometría del conjunto vertical. El rediseño se orientó a establecer una disposición común de sus guías y soportes, en lugar de limitar la intervención al contacto entre los carros y las guías.

Estos hallazgos aportaron el diagnóstico mecánico correspondiente al primer objetivo específico. La selección de componentes y el rediseño se desarrollaron en la segunda fase; la fabricación y la comprobación del mecanismo modificado correspondieron a la tercera. La inspección inicial no permitió cuantificar la precisión ni la repetibilidad del sistema original.

\section{Fase 2. Selección de componentes y rediseño mecánico}

La conservación de la arquitectura cartesiana respondió a que las deficiencias identificadas se concentraron en el guiado y la traslación. Por ello, se retomaron los diseños de los soportes de los ejes y de los motores de X y Y y se adaptaron a guías de acero de 12~mm. Esta decisión conservó parte de la geometría aprovechable del sistema, sin reutilizar las piezas impresas originales.

La selección de rodamientos lineales con carcasa facilitó la definición de uniones atornilladas entre los rodamientos y los carros. El modelado a partir de sus dimensiones estableció los espacios de montaje, el paso del eje perpendicular y el recorrido de la faja. En esta fase, el resultado correspondió a una configuración geométrica; su funcionamiento físico se comprobó después de la fabricación.

El conjunto central requirió un diseño específico porque articuló los desplazamientos en X y Y con el montaje del mecanismo vertical. La disposición de los rodamientos definió la separación entre guías y proporcionó una referencia común para los soportes de Z. La ubicación del efector final en la intersección de los ejes respondió a esa organización geométrica, sin constituir por sí sola una validación del posicionamiento del brazo.

Se conservaron las fajas dentadas porque estuvieron disponibles en el laboratorio. En Z se mantuvo el principio de transmisión mediante tornillo de avance y se modificaron los elementos de soporte para atender la desalineación observada. La intervención se concentró, por tanto, en integrar los componentes disponibles con el nuevo sistema de guiado.

La segunda fase aportó el diseño de los sistemas de traslación previsto en el primer objetivo específico. El ensamblaje CAD estableció la disposición de los componentes, pero no sustituyó la comprobación de las tolerancias de fabricación, la retención de los rodamientos ni el comportamiento del mecanismo ensamblado.

\section{Fase 3. Implementación mecánica}
\label{sec:discusion_fase3}

La fabricación y el ensamblaje evidenciaron diferencias entre la definición geométrica del modelo y las condiciones del montaje físico. Las dificultades se concentraron en la retención de los rodamientos del conjunto central, la alineación de las guías y el tensado de las fajas. Su resolución requirió iteraciones sobre los componentes y ajustes durante el montaje.

El PLA se utilizó por tratarse de un material común para impresión tridimensional. La impresión permitió fabricar nuevamente los componentes y modificar sus geometrías durante el desarrollo. Las tolerancias variaron entre piezas, y se identificó la posible necesidad de lijado en superficies curvas para mejorar el acoplamiento. Esta limitación mostró la importancia de comprobar físicamente los ajustes; no se documentó un valor único de tolerancia aplicable al conjunto.

Las tres iteraciones del conjunto central mostraron que facilitar la impresión no resolvió automáticamente la retención de los rodamientos. La división en tres piezas atendió la dificultad de fabricación, pero la disposición de los tornillos de presión se interpretó como una fuente de acciones opuestas que dificultaron la sujeción. La unión atornillada entre las secciones mantuvo los rodamientos en sus alojamientos durante las comprobaciones realizadas. El resultado vinculó la facilidad de fabricación con la necesidad de definir cómo quedaron unidas las piezas.

La alineación de las guías también condicionó el ensamblaje. El traslado de las distancias del modelo de Inventor al montaje proporcionó referencias para posicionar los soportes y facilitó esta operación. La ausencia de atascos en las comprobaciones posteriores fue compatible con una mejora del guiado, aunque no permitió separar el efecto de los nuevos rodamientos del efecto de los ajustes de alineación.

La fijación de las fajas mediante tornillos resolvió el montaje y el tensado necesarios para las comprobaciones iniciales. Sin embargo, no se midió la tensión ni se evaluó la durabilidad de las zonas perforadas. Del mismo modo, las arandelas se incorporaron con la intención de reducir el riesgo de aflojamiento, sin una evaluación específica de la conservación del apriete durante el uso prolongado. El alcance de estas soluciones quedó limitado a las condiciones comprobadas durante el montaje.

El soporte inferior desmontable de Z atendió la necesidad de retirar la garra para almacenar el equipo. El aporte del proyecto correspondió a su soporte y a la integración con el mecanismo vertical; la geometría original de la garra procedió de un modelo externo.\footnote{La identificación del autor y de la fuente del modelo de la garra quedó pendiente de documentación.}

La menor resistencia percibida, la ausencia de atascos durante las comprobaciones y la reducción visual del avance desigual del pórtico respaldaron una mejora cualitativa del movimiento. No se cuantificaron la fricción, la deflexión ni el desfase entre los extremos, por lo que estas observaciones no demostraron su eliminación completa ni una mejora medida de precisión o repetibilidad.

La tercera fase aportó la materialización y la comprobación mecánica preliminar del rediseño asociado con el primer objetivo específico. El desplazamiento a mano de los tres ejes estableció una base para las pruebas posteriores de accionamiento electrónico. Esta comprobación se distinguió del modo manual mediante mando inalámbrico y no sustituyó la validación experimental prevista en las fases posteriores.

\section{Fase 4. Accionamiento electrónico, operación manual y referenciación}
\label{sec:discusion_fase4}

La implementación progresiva separó el desarrollo del accionamiento, la lectura de los límites y la integración del mando. El movimiento por comandos proporcionó una base para incorporar posteriormente la interacción inalámbrica. Esta secuencia correspondió al orden de desarrollo descrito por el autor y no al orden de arranque del firmware integrado, que también incorporó funciones de etapas posteriores.

La distribución entre placas se describió según sus responsabilidades. La ESP32 concentró la interacción con el usuario, la OLED y los servomotores; la Portenta concentró el accionamiento de los ejes, los límites, la calibración y los estados generales. Aunque las órdenes del operador se originaron en la ESP32, las transacciones I2C fueron iniciadas por la Portenta. Por ello, la procedencia de una orden de movimiento no equivalió al papel de maestro del bus ni a la coordinación general del programa.

La separación entre el bus de la pantalla y el enlace de control permitió organizar ambos intercambios de manera independiente. Además, el firmware incorporó reintentos de inicialización de la OLED sin detener deliberadamente la atención del resto de las funciones. El intercambio de estados permitió que la interfaz presentara las direcciones ejecutadas y los límites informados por la Portenta, en lugar de limitarse a mostrar la entrada del joystick. No obstante, esos estados no constituyeron una medición independiente de la posición física del efector final.

El formato fijo de los paquetes estableció una representación común de las órdenes y de los estados en ambos programas. El CRC aportó una comprobación de integridad, mientras que el número de secuencia y la sesión de arranque permitieron distinguir la falta de actualización y el reinicio de la ESP32. Estas comprobaciones atendieron condiciones diferentes: la recepción de un mensaje con CRC válido no aseguró, por sí sola, que sus datos se hubieran actualizado recientemente. Del mismo modo, la presencia de CRC no garantizó la detección de cualquier alteración posible.

La entrega de una copia previamente preparada y la validación diferida de los mensajes recibidos separaron el procesamiento de datos de las funciones inmediatas de atención I2C. La prioridad asignada al envío de estados favoreció la actualización de la interfaz dentro de la planificación implementada, aunque pudo posponer un sondeo de control hasta la siguiente iteración. Sin mediciones del enlace, estos mecanismos no demostraron una latencia determinada ni ausencia de pérdidas de comunicación.

La zona neutra del mando delimitó la región en la que no se ordenó movimiento. Fuera de ella se seleccionó una dirección discreta, por lo que la amplitud del joystick no estableció una velocidad proporcional de traslación. Esta implementación atendió el movimiento manual por eje y aportó las funciones previstas en el segundo objetivo específico, sin demostrar por sí sola una mejora cuantificada del mantenimiento o del posicionamiento.

La doble aproximación a los finales estableció una secuencia definida de referenciación. El conteo entre extremos y la posterior ubicación en el centro proporcionaron un origen común para expresar los desplazamientos. Sin embargo, el conteo de pulsos generados representó la posición calculada por el programa y no una medición continua independiente del desplazamiento real del brazo.

La función denominada cinemática inversa correspondió a una conversión cartesiana de X/Y a destinos de pasos mediante factores de escala. Su alcance quedó limitado por los ejes efectivamente accionados: la recepción de un argumento Z no implicó posicionamiento cartesiano tridimensional. Asimismo, el movimiento por joystick no utilizó esa conversión, por lo que no se presentó como una ejecución de cinemática inversa mediante el mando.

La confirmación del recorrido físico de Y en 336~mm estableció la correspondencia entre el dato mecánico y la longitud utilizada para calcular el factor de pasos por milímetro. Esta correspondencia no constituyó por sí sola una validación de exactitud: la posición calculada también dependió del conteo entre extremos y de que el desplazamiento físico correspondiera a los pulsos generados. No se atribuyeron valores de precisión o repetibilidad a la rutina sin registros de posiciones obtenidas en ensayos.

Las comprobaciones de límites, la neutralización de órdenes y la validación de paquetes incorporaron respuestas de software frente a determinadas condiciones de fallo. Su presencia en el código no demostró una certificación de seguridad ni confirmó el comportamiento físico de todas las conexiones. La validación de las protecciones y del posicionamiento requirió distinguir los resultados observados en el prototipo de las funciones previstas en el firmware.
